\chapter[Calculus refresher]{Calculus refresher: tools for probability}
\label{app:calculus}
\chaptertag{A second encounter with limits, derivatives, integrals, and series}

Calculus returns in probability whenever a sum becomes an area, a product
needs an approximation, or a parameter must be chosen to maximize a
likelihood. This appendix rebuilds those tools through worked examples.
Read with a pencil: pause before each calculation and predict its sign,
scale, or limiting value. The exercises at the end have worked answers.

\begin{bigidea}[title={Choose a route through the refresher}]
For birthdays, review logarithms, finite sums, and Taylor approximation,
then read Section~\ref{sec:calc-birthdays}. For derangements, review
series and the exponential function, then read
Section~\ref{sec:calc-derangements}. For continuous random variables,
start with definite integrals and the fundamental theorem in
Section~\ref{sec:calc-integrals}. The final beta-integral example is an
optional bridge to Bayesian inference.
\end{bigidea}

\section{Algebra, functions, and finite sums}

\subsection*{Keep track of the variable}
In $f(x)=\lambda e^{-\lambda x}$, $x$ is the input and $\lambda>0$ is a
fixed parameter. Differentiating with respect to $x$ holds $\lambda$
constant. Later, maximizing a likelihood in $\lambda$ holds the observed
data constant instead. Before differentiating, say which symbol is allowed
to vary.

The natural logarithm $\log x$ is defined for $x>0$ and is the inverse
of the exponential function. Its essential rules are
\[
 e^{u+v}=e^u e^v,\qquad
 \log(ab)=\log a+\log b,\qquad
 \log(a^r)=r\log a\quad(a,b>0).
\]
Also $e^0=1$, $\log1=0$, and $e^{\log a}=a$.
Taking logarithms turns a product of positive factors into a sum.
It does not turn a sum into a sum of logarithms:
$\log(1+1)=\log2\ne0=\log1+\log1$.

\subsection*{A sum is an instruction, not a new operation}
The index in $\sum_{j=0}^{n-1}j$ runs through $0,1,\ldots,n-1$;
there are $n$ terms. The two sums needed for birthday estimates are
\[
 \sum_{j=0}^{n-1}j=\frac{n(n-1)}2,\qquad
 \sum_{j=0}^{n-1}j^2=\frac{n(n-1)(2n-1)}6.
\]
Pair the first sum with its reversed order: each of the $n$ paired
terms is $n-1$. For the second identity, induction works: increasing
$n$ to $n+1$ adds $n^2$ to both sides. Both sides start at zero
when $n=1$.

The product notation $\prod_{j=0}^{n-1}(1-j/d)$ means
$1(1-1/d)\cdots(1-(n-1)/d)$. A sum with no terms is zero; a product
with no factors is one. These conventions make the cases $n=0$ work
without separate algebraic exceptions.

Factorials are finite products, $n!=1\cdot2\cdots n$, with $0!=1$.
Cancel common factors before using a calculator:
\[
 \frac{n!}{(n-2)!}=n(n-1),\qquad
 \frac{\binom nk(n-k)!}{n!}=\frac1{k!}.
\]
The second cancellation explains the simple denominators in the
derangement probability.

\section{Limits and continuity}

The statement $f(x)\to L$ as $x\to a$ says that $f(x)$ can be made
as close as desired to $L$ by taking $x$ sufficiently close to $a$,
with $x\ne a$. The value at $a$ itself is a separate question.
For example,
\[
 \lim_{x\to1}\frac{x^2-1}{x-1}
 =\lim_{x\to1}(x+1)=2.
\]
The fraction is undefined at $1$, but the nearby values have a limit.
A function is \emph{continuous} at $a$ when the limit equals $f(a)$.
Polynomials are continuous everywhere; rational functions are continuous
where their denominators are nonzero. Exponentials are continuous, and
logarithms are continuous on their positive domain.

For a sequence $a_n$, the notation $a_n\to L$ means that all sufficiently
late terms stay arbitrarily close to $L$. Thus $1/n\to0$, while
$(-1)^n$ has no limit. The expression
\[
 \left(1+\frac1n\right)^n\longrightarrow e\approx2.7182818
\]
is one familiar appearance of the exponential constant. Another,
$e=\sum_{k=0}^{\infty}1/k!$, will follow from Taylor's theorem below.

\begin{warningbox}
A small term can accumulate. Although $1/n\to0$, a sum of $n$ copies
of $1/n$ is always one. Likewise, replacing every birthday factor by
something ``close'' to it needs an estimate for the total error.
\end{warningbox}

\subsection*{Two practical limit tools}
If $0\le a_n\le b_n$ and $b_n\to0$, then $a_n\to0$ by squeezing.
For example, $0\le n/(n^2+1)\le1/n$.
For quotients with the indeterminate forms $0/0$ or $\infty/\infty$,
l'H\^opital's rule can replace a quotient by the quotient of its
derivatives, provided the functions are differentiable nearby, the
denominator derivative is nonzero there, and the latter quotient has a
limit. It is not a rule for arbitrary fractions. A useful example is
\[
 \lim_{x\to\infty}\frac{x}{e^x}
 =\lim_{x\to\infty}\frac1{e^x}=0.
\]
Repeated application shows $x^m e^{-\lambda x}\to0$ for a fixed
nonnegative integer $m$ and $\lambda>0$. This justifies the boundary
terms in exponential moment calculations.

\section{Derivatives describe local change}
\label{sec:calc-derivatives}

The derivative is the limiting slope
\[
 f'(a)=\lim_{h\to0}\frac{f(a+h)-f(a)}h.
\]
For $f(x)=x^2$, the difference quotient is $2a+h$, so $f'(a)=2a$.
For small $h$, the tangent-line estimate is
$f(a+h)\approx f(a)+f'(a)h$. This is an approximation near $a$;
Taylor's theorem will quantify its error.

\begin{center}
\begin{tabular}{@{}ll@{}}
\toprule
Function or operation & Derivative \\\midrule
$c$; $x^r$ & $0$; $r x^{r-1}$ where defined \\
$e^x$; $\log x$ & $e^x$; $1/x$ for $x>0$ \\
$af+bg$ & $af'+bg'$ \\
$fg$ & $f'g+fg'$ \\
$f/g$ & $(f'g-fg')/g^2$ where $g\ne0$ \\
$f(g(x))$ & $f'(g(x))g'(x)$ \\\bottomrule
\end{tabular}
\end{center}

The last identity is the \emph{chain rule}: differentiate the outer
function and multiply by the derivative of the inner function. For example,
\[
 \frac{d}{dx}e^{-3x}=-3e^{-3x},\qquad
 \frac{d}{dx}\log(1-x)=-\frac1{1-x}\quad(x<1).
\]
The product and chain rules together give
$\frac{d}{dx}(x^2e^{-3x})=(2x-3x^2)e^{-3x}$.
The second derivative $f''$ describes change in the slope: $f''>0$
means the graph bends upward, and $f''<0$ means it bends downward.

\begin{example}[Differentiation remembers a counting identity]
Differentiate the \emph{finite} binomial identity
$(1+x)^n=\sum_{k=0}^n\binom nkx^k$. For $n\ge1$,
\[
 n(1+x)^{n-1}=\sum_{k=1}^n k\binom nkx^{k-1}.
\]
At $x=1$ this is $\sum_{k=0}^n k\binom nk=n2^{n-1}$.
The same quantity counts subsets with a distinguished member: choose
the member in $n$ ways, then decide which of the other $n-1$ objects to include.
\end{example}

\section{Optimization: logarithms simplify a likelihood}

A differentiable function can have an interior maximum only where its
derivative is zero. That is a candidate, not a conclusion: check the
sign of the derivative or the second derivative, and check boundaries.

Suppose $s$ successes and $f=n-s$ failures have been observed in
independent trials with unknown success probability $p$. For a specified
observed string the likelihood is $L(p)=p^s(1-p)^f$. If only the count
is recorded, the extra factor $\binom ns$ does not depend on $p$ and
does not change the maximizing value.

When $0<s<n$, maximize its logarithm on $0<p<1$:
\[
 \ell(p)=s\log p+f\log(1-p),\qquad
 \ell'(p)=\frac{s}{p}-\frac{f}{1-p}.
\]
Logarithm is increasing, so the maximizer is preserved. Setting the
derivative to zero gives $s(1-p)=fp$, hence $p=s/n$. Moreover,
\[
 \ell''(p)=-\frac{s}{p^2}-\frac{f}{(1-p)^2}<0,
\]
so this is the unique interior maximum. The likelihood is zero at
both endpoints when both counts are positive.
If $s=0<n$, $L(p)=(1-p)^n$ is maximized at $p=0$; if $s=n>0$,
$L(p)=p^n$ is maximized at $p=1$. With no observations, the likelihood
is constant and selects no unique parameter.

\begin{checkpoint}
With 7 successes in 10 trials the maximizing parameter is $.7$.
Explain why this calculation treats the data as fixed and $p$ as
variable. A likelihood is not, by itself, a probability density for $p$.
\end{checkpoint}

\section{Integrals: area and accumulated change}
\label{sec:calc-integrals}

Divide $[a,b]$ into narrow intervals of widths $\Delta x_i$, and choose
a sample point $t_i$ in each one. The sum
$\sum_i f(t_i)\Delta x_i$ adds the areas of thin rectangles.
For a continuous $f$, these sums tend to the definite integral
$\int_a^b f(x)\,dx$ as the largest width tends to zero.
If $f\ge0$, this is area; below the axis, contributions are signed.
Probability densities must be nonnegative so that their areas are probabilities.

\begin{theorem}[Fundamental theorem of calculus]
If $f$ is continuous on $[a,b]$, then
$A(x)=\int_a^x f(t)\,dt$ satisfies $A'(x)=f(x)$ in the interior.
If $G'=f$, then
\[
 \int_a^b f(x)\,dx=G(b)-G(a).
\]
\end{theorem}
An antiderivative turns an area calculation into an endpoint calculation.
The indefinite integral $\int f(x)\,dx=G(x)+C$ is a family of functions;
the definite integral is a number. The arbitrary constant cancels between
the endpoints.

The power rule gives
\[
 \int x^r\,dx=\frac{x^{r+1}}{r+1}+C\quad(r\ne-1),\qquad
 \int\frac{dx}{x}=\log|x|+C\quad(x\ne0).
\]
Use intervals on which the integrand is defined. Integrals distribute over
finite sums and constant multiples, but not over products.

\begin{example}[Normalize, accumulate, average]
Suppose a proposed density is $cx$ on $[0,1]$ and zero elsewhere.
Normalization requires $1=\int_0^1 cx\,dx=c/2$, so $c=2$.
For $0\le t\le1$, its CDF is
$F(t)=\int_0^t2x\,dx=t^2$. Its first two moments are
\[
 \E[X]=\int_0^1 x(2x)\,dx=\frac23,\qquad
 \E[X^2]=\int_0^1 x^2(2x)\,dx=\frac12.
\]
Consequently $\Var(X)=1/2-4/9=1/18$.
The event $X\ge1/2$ has probability
$\int_{1/2}^1 2x\,dx=[x^2]_{1/2}^1=3/4$.
\end{example}

\begin{center}
\begin{tikzpicture}[x=3.7cm,y=1.65cm]
 \fill[mint!25] (.5,0)--(.5,1)--(1,2)--(1,0)--cycle;
 \draw[->,muted] (0,0)--(1.15,0) node[right] {$x$};
 \draw[->,muted] (0,0)--(0,2.3) node[above] {$f(x)$};
 \draw[indigo,very thick] (0,0)--(1,2);
 \node[below] at (.5,0) {$1/2$}; \node[below] at (1,0) {$1$};
 \node[left] at (0,2) {$2$};
 \node[fill=paper,inner sep=2pt] at (.78,.45) {area $3/4$};
 \begin{scope}[xshift=6cm,y=3.3cm]
  \draw[->,muted] (0,0)--(1.15,0) node[right] {$x$};
  \draw[->,muted] (0,0)--(0,1.15) node[above] {$F(x)$};
  \draw[violet,very thick,domain=0:1,samples=50] plot (\x,{\x*\x});
  \draw[dashed,muted] (.5,0)--(.5,.25)--(0,.25);
  \node[below] at (.5,0) {$1/2$}; \node[below] at (1,0) {$1$};
  \node[left] at (0,.25) {$1/4$}; \node[left] at (0,1) {$1$};
 \end{scope}
\end{tikzpicture}
\end{center}
The shaded area under the density equals $1-F(1/2)$.
The height of the CDF already represents accumulated probability.

\section{Substitution and integration by parts}
\label{sec:calc-techniques}

\subsection*{Substitution reverses the chain rule}
If $u=g(x)$, then $du=g'(x)\,dx$. For example,
\[
 \int_0^1 2x e^{-x^2}\,dx
 =\int_0^1 e^{-u}\,du=1-e^{-1},\qquad u=x^2.
\]
Both endpoints have been changed into $u$-values. In general,
\[
 \int_a^b h(g(x))g'(x)\,dx
 =\int_{g(a)}^{g(b)}h(u)\,du
\]
under the usual continuity and differentiability assumptions.
This identity also handles reversed endpoints correctly.

\begin{warningbox}
After a substitution, change both the differential and the bounds.
The expression $\int_a^b h(u)\,du$ still uses $u$ as its input; old
$x$-bounds cannot be retained unless they happen to give the same numbers.
\end{warningbox}

\subsection*{Integration by parts reverses the product rule}
Integrating $(uv)'=u'v+uv'$ gives
\[
 \int_a^b u(x)v'(x)\,dx
 =[u(x)v(x)]_a^b-\int_a^b u'(x)v(x)\,dx.
\]
Choose $u$ so that differentiating it simplifies the remaining integral.
For example, $u=x$, $v'=e^x$, and $v=e^x$ give
\[
 \int_0^1 xe^x\,dx=[xe^x]_0^1-\int_0^1e^x\,dx
 =e-(e-1)=1.
\]
This technique will also compute the mean and variance of an exponential
waiting time and the beta integral used in Bayesian updating.

\section{Improper integrals and infinite expectations}
\label{sec:calc-improper}

An infinite endpoint is a limit:
$\int_0^\infty f(x)\,dx=\lim_{R\to\infty}\int_0^R f(x)\,dx$,
when that limit exists. A singular finite endpoint is treated similarly;
for example $\int_0^1 x^{-1/2}\,dx=\lim_{\varepsilon\downarrow0}
[2\sqrt{x}]_{\varepsilon}^1=2$.
Do not substitute infinity as if it were an ordinary number.

\begin{example}[Exponential waiting times]
For $\lambda>0$, let $f(x)=\lambda e^{-\lambda x}$ on $x\ge0$.
Its integral on $[0,R]$ is $1-e^{-\lambda R}$, tending to one.
Thus it is a density, and
$\P(X>t)=\int_t^\infty\lambda e^{-\lambda x}\,dx=e^{-\lambda t}$
for $t\ge0$.
Integration by parts on a finite interval gives
\[
 \int_0^R x\lambda e^{-\lambda x}\,dx
 =-R e^{-\lambda R}+\frac{1-e^{-\lambda R}}{\lambda}.
\]
Since $R e^{-\lambda R}\to0$, the mean is $1/\lambda$.
Another integration by parts, using $R^2e^{-\lambda R}\to0$, gives
\[
 \E[X^2]=2\int_0^\infty xe^{-\lambda x}\,dx=\frac2{\lambda^2},
 \qquad \Var(X)=\frac1{\lambda^2}.
\]
The median solves $e^{-\lambda m}=1/2$, hence $m=\log2/\lambda$.
It differs from the mean because the long right tail pulls the mean upward.
\end{example}

\begin{example}[A density can have no finite mean]
Let $f(x)=1/(1+x)^2$ on $x\ge0$ and zero elsewhere.
Its total area is one, because
$\int_0^R(1+x)^{-2}\,dx=1-1/(1+R)\to1$.
But
\[
 \int_0^R\frac{x}{(1+x)^2}\,dx
 =\log(1+R)+\frac1{1+R}-1\longrightarrow\infty.
\]
The distribution is valid and its expected value is $+\infty$.
Normalizing a density and proving its mean finite are different calculations.
\end{example}

The general tail check is
$\int_1^\infty x^{-r}\,dx<\infty$ exactly when $r>1$.
Integrate $x^{-r}$ for $r\ne1$ and use $\log x$ for $r=1$ to verify
the boundary between convergence and divergence.

\section{Sequences, series, and geometric probabilities}
\label{sec:calc-series}

An infinite series $\sum_{k=0}^{\infty}a_k$ is defined as the limit of
its partial sums $S_m=\sum_{k=0}^m a_k$, when this limit exists.
The terms $a_k$ and their sums $S_m$ are different sequences.
Convergence requires $a_k\to0$, since $a_k=S_k-S_{k-1}$, but that
condition alone is insufficient. The harmonic series diverges:
\[
 1+\frac12+\left(\frac13+\frac14\right)
 +\left(\frac15+\cdots+\frac18\right)+\cdots.
\]
Every displayed group after the first two terms contributes at least $1/2$.
There are arbitrarily many such groups.

\subsection*{The geometric series}
For $r\ne1$, subtracting $r$ times a finite sum from itself gives
\[
 1+r+\cdots+r^m=\frac{1-r^{m+1}}{1-r}.
\]
If $|r|<1$, the final power tends to zero, so
\[
 \sum_{k=0}^{\infty}r^k=\frac1{1-r}.
\]
For $0\le r<1$, the tail after $r^m$ is exactly $r^{m+1}/(1-r)$.
The condition $|r|<1$ matters: substituting $r=2$ into the infinite
formula would produce a negative answer for a growing positive sum.

\begin{example}[Wait until the first success]
In independent trials with success probability $0<p\le1$, let $N$
count trials up to and including the first success. With $r=1-p$,
$\P(N=k)=pr^{k-1}$ for $k\ge1$, and the geometric series checks
that these probabilities sum to one. For $0<p<1$, differentiating
the geometric power series inside $|r|<1$ gives
\[
 \sum_{k=1}^{\infty}kr^{k-1}=\frac1{(1-r)^2},\qquad
 \E[N]=\sum_{k=1}^{\infty}kpr^{k-1}=\frac1p.
\]
Power series can be differentiated term by term strictly inside their
interval of convergence. Alternatively, differentiate the finite
geometric identity first and let its length tend to infinity;
the extra terms involve $m r^m$ and tend to zero.
For $p=1$, $N=1$ with certainty, so the same mean formula holds directly.
\end{example}

\subsection*{Alternating sums come with a useful error bound}
If $b_k$ decreases to zero through positive values, the partial sums of
$b_0-b_1+b_2-\cdots$ alternate above and below its limit.
The magnitude of the error is at most the first omitted term. To see
why, group consecutive terms: the even partial sums decrease, the odd
partial sums increase, and their difference tends to zero. The limit is
trapped between the two adjacent partial sums.
For a strictly decreasing positive tail, the limit is strictly between them.
This is exactly the estimate needed for derangements.

\section{Taylor polynomials and a reason to trust them}
\label{sec:calc-taylor}

The degree-$m$ Taylor polynomial at $a$ matches the first $m$ derivatives
of a function there:
\[
 T_m(x)=\sum_{k=0}^m\frac{f^{(k)}(a)}{k!}(x-a)^k.
\]
The factorial appears because the $k$th derivative of $(x-a)^k$ at
$a$ is $k!$. The first two terms give the tangent line; the quadratic
term also records curvature.

\begin{theorem}[Taylor's error estimate]
Suppose $f$ has continuous derivatives through order $m+1$ on the
interval between $a$ and $x$, and $|f^{(m+1)}(t)|\le M$ there.
Then
\[
 |f(x)-T_m(x)|\le\frac{M|x-a|^{m+1}}{(m+1)!}.
\]
\end{theorem}
One way to understand the bound is to repeatedly integrate the derivative
error using the fundamental theorem. The remainder has the form
\[
 R_m(x)=\frac1{m!}\int_a^x (x-t)^m f^{(m+1)}(t)\,dt.
\]
Taking absolute values and integrating $M|x-t|^m/m!$ gives the estimate,
with the same conclusion when $x<a$ after reversing the limits.

\subsection*{The exponential series}
All derivatives of $e^x$ are $e^x$. Expanding at zero gives
\[
 e^x=1+x+\frac{x^2}{2!}+\cdots+\frac{x^m}{m!}+R_m(x),
 \qquad |R_m(x)|\le\frac{e^{|x|}|x|^{m+1}}{(m+1)!}.
\]
For fixed $x$, the successive factorial-denominator terms have ratio
$|x|/(m+2)$, eventually less than $1/2$. They therefore tend to zero,
proving
\[
 e^x=\sum_{k=0}^{\infty}\frac{x^k}{k!}\quad\text{for every real }x.
\]
For example, $e^{.1}\approx1+.1+.1^2/2=1.105$; the error is at most
$e^{.1}(.1)^3/6<.000185$. The estimate is accurate because it includes
a small remainder bound, not merely because a few terms were retained.

\subsection*{Logarithms from a geometric series}
For $0\le x<1$, integrate the finite geometric identity for $1/(1-t)$
from $0$ to $x$. For $m\ge1$ this gives
\[
 -\log(1-x)=\sum_{k=1}^m\frac{x^k}{k}
       +\int_0^x\frac{t^m}{1-t}\,dt.
\]
Since $1/(1-t)\le1/(1-x)$ on this interval, the remainder is between
zero and $x^{m+1}/[(m+1)(1-x)]$. It tends to zero for fixed $x<1$,
so
\[
 \log(1-x)=-x-\frac{x^2}{2}-\frac{x^3}{3}-\cdots.
\]
At $x=.1$, the first two terms give $-.105$, with error less than
$.1^3/[3(.9)]<.000371$. At $x=1$, the logarithm is undefined and
the corresponding harmonic series diverges. An expansion's domain
is part of the formula.

\begin{warningbox}
A function is not automatically equal to its infinite Taylor series
just because many derivatives exist. For the exponential and logarithm
used here, the explicit remainders above prove equality on the stated
domains. For an approximation, keep enough terms to control the error
at the inputs you actually use.
\end{warningbox}

\section{Birthdays: logarithms turn the product into a sum}
\label{sec:calc-birthdays}

In the independent uniform $d$-day birthday model, write $q_d(n)$ for
the chance of a collision and $r_d(n)=1-q_d(n)$ for the chance of no
collision. For $1\le n\le d$,
\[
 r_d(n)=\prod_{j=0}^{n-1}(1-j/d),\qquad
 \log r_d(n)=\sum_{j=0}^{n-1}\log(1-j/d).
\]
Keeping just the linear term $\log(1-x)\approx-x$ yields
\[
 \log r_d(n)\approx-\frac{n(n-1)}{2d},\qquad
 q_d(n)\approx1-\exp\!\left(-\frac{n(n-1)}{2d}\right).
\]
There are $\binom n2$ pairs, so this exponent already suggests why a
class much smaller than $d$ can produce a match.

\subsection*{An error estimate, with a direction}
For $0\le x\le c<1$,
\[
 -\log(1-x)-x=\int_0^x\frac{t}{1-t}\,dt.
\]
Bounding the integrand between $t$ and $t/(1-c)$ gives
\[
 -x-\frac{x^2}{2(1-c)}\le\log(1-x)
 \le-x-\frac{x^2}{2}\le-x.
\]
Set
\[
 c=\frac{n-1}{d},\qquad S_1=\frac{n(n-1)}{2d},\qquad
 S_2=\frac{n(n-1)(2n-1)}{6d^2}.
\]
Sum the inequalities over the birthday factors and exponentiate:
\[
 \boxed{\quad
 1-e^{-S_1-S_2/2}\ \le\ q_d(n)\ \le\
 1-e^{-S_1-S_2/[2(1-c)]}.
 \quad}
\]
In particular, the simpler estimate $1-e^{-S_1}$ is a lower bound.
Its omitted correction in the logarithm is at most $S_2/[2(1-c)]$.
The approximation is useful when this correction is small; having
each individual factor close to one is not by itself an error guarantee.
For $n=0$, the exact collision probability is zero; for $n>d$, it is one,
and we use those exact cases instead of the logarithmic approximation.

\begin{center}
\small
\begin{tabular}{@{}rrrr@{}}
\toprule
$n$ on 365 days & $1-e^{-S_1}$ & $1-e^{-S_1-S_2/2}$ & Exact $q_{365}(n)$\\\midrule
10&.115991&.116936&.116948\\
22&.468938&.475496&.475695\\
23&.500002&.507073&.507297\\
30&.696320&.705916&.706316\\
50&.965131&.970040&.970374\\\bottomrule
\end{tabular}
\end{center}
The entries are rounded, not exact inequalities between the printed
decimals. At $n=23$, the rigorous bounds give
$.50707<q_{365}(23)<.50753$; direct multiplication gives about $.507297$.

\subsection*{Where the number 23 comes from}
Set the first approximation equal to $1/2$. Taking logarithms gives
$n(n-1)=2d\log2$, hence the positive solution
\[
 n\approx\frac{1+\sqrt{1+8d\log2}}2.
\]
For $d=365$, this is about $22.99994$, suggesting a threshold of 23.
The exact probabilities at 22 and 23, $.475695$ and $.507297$,
verify the conclusion. This use of a smooth expression estimates an
integer class size; the exact model still checks the adjacent integers.

\begin{labbox}[title={Calculate with small tails intact}]
For numerical work, sum $\log(1-j/d)$ and then exponentiate instead of
multiplying many factors. Software functions named \texttt{log1p(z)}
compute $\log(1+z)$ accurately near zero, and \texttt{expm1(z)} computes
$e^z-1$. Thus $q_d(n)=-\texttt{expm1}(\sum_j\texttt{log1p}(-j/d))$
is a useful evaluation for $n\le d$. It evaluates the exact product
formula with floating-point arithmetic; it is not the Taylor approximation.
\end{labbox}

\section{Derangements: why the limiting probability is one over e}
\label{sec:calc-derangements}

Section~\ref{sec:derangements} obtained an exact finite expression:
\[
 \frac{D_n}{n!}=1-1+\frac1{2!}-\frac1{3!}+\cdots+\frac{(-1)^n}{n!}.
\]
Now set $x=-1$ in the exponential series. Its sum is $e^{-1}$.
The positive terms in every omitted tail strictly decrease to zero, so
the alternating-series estimate gives, for $n\ge0$,
\[
 \left|\frac{D_n}{n!}-e^{-1}\right|<\frac1{(n+1)!}.
\]
Even partial sums are above the limit and odd partial sums below it.
For example,
\[
 \frac{44}{120}<e^{-1}<\frac{265}{720},
 \qquad e^{-1}\approx.3678794412.
\]
At $n=10$, the exact probability is
$1334961/3628800\approx.3678794643$, within $1/11!<2.51\cdot10^{-8}$
of the limit. This is why the approximation becomes excellent so quickly.

\begin{proposition}[Recover an integer count by rounding]
For every $n\ge1$, $D_n$ is the nearest integer to $n!/e$.
\end{proposition}
\begin{proof}
Multiply the strict error bound by $n!$ to obtain
$|D_n-n!/e|<1/(n+1)\le1/2$. Thus the integer $D_n$ lies strictly
within half a unit of $n!/e$.
\end{proof}
The condition $n\ge1$ matters: $D_0=1$, whereas $1/e$ rounds to zero.
For large $n$, the exact recurrence is preferable for computer
calculations because ordinary floating-point arithmetic may not retain
unit accuracy in $n!/e$.

\subsection*{A further connection: the number of fixed points}
For a uniform permutation, the probability of exactly $k$ fixed points is
\[
 \frac{\binom nk D_{n-k}}{n!}
 =\frac1{k!}\frac{D_{n-k}}{(n-k)!}\longrightarrow\frac{e^{-1}}{k!}
\]
as $n\to\infty$ with $k$ fixed. These limiting probabilities sum to one
by the exponential series. They form the Poisson distribution with mean
one, introduced in Chapter 5. Already, indicators tell us that the
expected number of fixed points is exactly one for every $n\ge1$.

Birthdays and derangements both lead to exponentials, through different
calculations: birthdays use a logarithmic expansion of a product;
derangements produce a finite alternating exponential sum by
inclusion--exclusion. In each case, the remainder explains the accuracy.

\section{Optional: a beta integral for Bayesian updating}
\label{sec:calc-beta}

For nonnegative integers $s,f$, consider
\[
 I_{s,f}=\int_0^1p^s(1-p)^f\,dp.
\]
This integral normalizes the likelihood from $s$ successes and $f$
failures when a uniform prior density is used. For $s\ge1$, integration
by parts with $u=p^s$ and $v'= (1-p)^f$ gives
\[
 I_{s,f}=\frac{s}{f+1}I_{s-1,f+1}.
\]
The boundary term is zero at both ends. Starting from
$I_{0,m}=1/(m+1)$ and applying the recurrence $s$ times gives
\[
 I_{s,f}=\frac{s!f!}{(s+f+1)!}.
\]
For example, $I_{3,1}=1/20$, so $20p^3(1-p)$ is a density on $[0,1]$.
Its mean is $I_{4,1}/I_{3,1}=(1/30)/(1/20)=2/3$.
More generally the mean is
$I_{s+1,f}/I_{s,f}=(s+1)/(s+f+2)$.
Chapter 9 interprets this as a posterior mean and predictive probability.

\clearpage
\section{Practice: rebuild the tools}
\label{sec:calc-exercises}

\begin{exercisebox}
\begin{enumerate}
 \item Simplify $\binom n3(n-3)!/n!$ for $n\ge3$. Evaluate
 $\sum_{j=0}^{5}j$ and $\sum_{j=0}^{5}j^2$. \exlevel{Fluency}
 \item Find $\lim_{x\to2}(x^2-4)/(x-2)$ and
 $\lim_{n\to\infty}n/(n^2+1)$. Explain why direct substitution is
 insufficient for the first. \exlevel{Fluency}
 \item Differentiate $x^2e^{-x}$ and $\log(1-x^2)$, specifying
 the real domain of the second function. \exlevel{Fluency}
 \item Use a tangent line at zero to approximate $e^{.05}$.
 Use Taylor's estimate to bound its absolute error. \exlevel{Reasoning}
 \item Maximize $p^8(1-p)^2$ on $[0,1]$. Explain why maximizing
 $p^{10}$ needs a boundary check instead. \exlevel{Fluency}
 \item Differentiate the binomial identity twice and set $x=1$
 to prove $\sum_{k=0}^n k(k-1)\binom nk=n(n-1)2^{n-2}$ for $n\ge2$.
 \exlevel{Reasoning}
 \item Let a density be $cx^2$ on $[0,1]$, zero elsewhere.
 Find $c$, its CDF, $\P(X>1/2)$, its mean, and its variance.
 \exlevel{Fluency}
 \item Evaluate $\int_0^2 2x e^{-x^2}\,dx$ using substitution.
 Show the new bounds. \exlevel{Fluency}
 \item Evaluate $\int_0^1 x e^{2x}\,dx$ by parts.
 \exlevel{Fluency}
 \item Normalize $c e^{-3x}$ on $x\ge0$. Find its mean, variance,
 and median. \exlevel{Fluency}
 \item For density $1/(1+x)^2$ on $x\ge0$, find the survival
 probability $\P(X>t)$ for $t\ge0$ and the median. Explain how a
 finite median is compatible with an infinite mean. \exlevel{Reasoning}
 \item Sum $\sum_{k=4}^{\infty}(1/2)^k$. If a geometric waiting
 time has success probability $.2$, find its mean and
 $\P(N>5)$. \exlevel{Fluency}
 \item Explain why terms tending to zero do not guarantee that a
 series converges. Contrast $\sum 1/k$ with $\sum (1/2)^k$.
 \exlevel{Reasoning}
 \item Approximate $\log(.95)$ using $-x-x^2/2$ at $x=.05$.
 Bound the error using the logarithmic remainder. \exlevel{Fluency}
 \item For $n=23,d=365$, compute $S_1,S_2,c$ and use the two
 logarithmic bounds to bracket the collision probability. Compare
 with the first-order estimate. \exlevel{Reasoning}
 \item Compute $D_8$ using the exact recurrence. Bound the error
 in $D_8/8!\approx e^{-1}$, determine which is larger, and check
 the nearest-integer conclusion. \exlevel{Reasoning}
 \item Evaluate $\int_0^1p^2(1-p)^2\,dp$ both by expanding and
 by the beta formula. Normalize the integrand and find the resulting
 density's mean. \exlevel{Fluency}
 \item For a uniform permutation of four objects, find the probability
 of exactly two fixed points. Compare it with the limiting value
 $e^{-1}/2!$. Explain why a limiting formula need not be accurate
 for such a small permutation. \exlevel{Reasoning}
\end{enumerate}
\end{exercisebox}

\section{Worked answers to the refresher}
\begin{enumerate}
 \item Cancel factorials to get $1/3!=1/6$. The sums are
 $5\cdot6/2=15$ and $5\cdot6\cdot11/6=55$.
 \item Away from $x=2$, the fraction equals $x+2$, so the first
 limit is $4$. At $2$ the original denominator is zero.
 For the sequence, $0\le n/(n^2+1)\le1/n\to0$.
 \item The product rule gives $(2x-x^2)e^{-x}$.
 The logarithm requires $1-x^2>0$, so $-1<x<1$;
 its derivative is $-2x/(1-x^2)$ by the chain rule.
 \item The approximation is $1+.05=1.05$. On $[0,.05]$,
 the second derivative of $e^x$ is at most $e^{.05}$, so the error
 is at most $e^{.05}(.05)^2/2<.001315$. The positive quadratic
 and later terms show the tangent estimate is too small.
 \item On $0<p<1$, the log derivative is $8/p-2/(1-p)$,
 which vanishes at $p=.8$. Its derivative is negative, so this is
 the maximum; the original function is zero at both endpoints.
 The function $p^{10}$ increases on $[0,1]$ and is maximized at $1$.
 \item Twice differentiating gives
 $n(n-1)(1+x)^{n-2}=\sum_{k=2}^n k(k-1)\binom nkx^{k-2}$.
 Set $x=1$. The terms $k=0,1$ contribute zero to the requested sum.
 \item Normalization gives $c/3=1$, hence $c=3$.
 The CDF is $0$ below zero, $x^3$ on $[0,1]$, and $1$ above one.
 Thus $\P(X>1/2)=7/8$,
 $\E[X]=3\int_0^1x^3\,dx=3/4$, and
 $\E[X^2]=3/5$. The variance is $3/5-9/16=3/80$.
 \item With $u=x^2$, $du=2x\,dx$ and the bounds become $0,4$.
 The integral is $\int_0^4e^{-u}\,du=1-e^{-4}$.
 \item Set $u=x$ and $v'=e^{2x}$, so $v=e^{2x}/2$.
 The result is
 $[xe^{2x}/2]_0^1-\frac12\int_0^1e^{2x}\,dx
 =e^2/2-(e^2-1)/4=(e^2+1)/4$.
 \item The total integral is $c/3$, so $c=3$.
 This is an exponential density with rate $3$: its mean is $1/3$,
 variance $1/9$, and median $(\log2)/3$.
 \newpage
 \item Integrating the tail gives
 $\P(X>t)=1/(1+t)$. Setting this to $1/2$ gives median $1$.
 A median divides probability in half; a mean weights outcomes by
 their magnitudes. The tail is heavy enough that those weighted
 contributions have infinite total.
 \item The geometric tail is $(1/2)^4/(1-1/2)=1/8$.
 The waiting-time mean is $1/.2=5$. The event $N>5$ means five
 initial failures, with probability $.8^5=.32768$.
 \item Harmonic terms tend to zero, but grouping indices between
 consecutive powers of two supplies infinitely many contributions
 of at least $1/2$. The geometric partial sums instead approach a
 finite limit because the remainder is a vanishing power divided
 by $1-r$.
 \item The approximation is $-.05-.05^2/2=-.05125$.
 The omitted magnitude is at most
 $.05^3/[3(.95)]<.00004386$. The exact logarithm is more negative,
 so $-.05129386<\log(.95)<-.05125$.
 \item Here $S_1=253/365$, $S_2=3795/133225$, and $c=22/365$.
 The lower and upper expressions are approximately $.507072664$
 and $.507522764$. Thus an outward-rounded bracket is
 $.50707<q_{365}(23)<.50753$. The first-order estimate is
 $.500001752$, which is noticeably lower.
 \item Continuing from $D_6=265$ gives $D_7=7\cdot265-1=1854$
 and $D_8=8\cdot1854+1=14833$.
 Since $8$ is even, $D_8/8!>e^{-1}$, with error less than
 $1/9!\approx.000002756$. Multiplying by $8!$ puts $14833$
 within $1/9$ of $8!/e$, so it is the nearest integer.
 \item Expand $p^2(1-p)^2=p^2-2p^3+p^4$; its integral is
 $1/3-1/2+1/5=1/30$, also $2!2!/5!$.
 The density is $30p^2(1-p)^2$, with mean
 $30(1/4-2/5+1/6)=1/2$.
 \item Choose the fixed pair and swap the other two, giving
 $\binom42D_2/4!=6/24=1/4$.
 The limiting probability is about $.18394$. In the exact formula
 $\frac1{2!}D_2/2!$, only two objects remain to derange; their
 probability $1/2$ is still far from $1/e$.
\end{enumerate}
